\section{Serving implications}
\label{sec:serving}

\ingram{} changes the state that a served model stores for each session, and where that state can
reside. We first report measured counts and then present two scenarios from a roofline model. The
scenarios assume a model with fewer attention layers; Section~\ref{sec:results} provides the
evidence, at up to 115M parameters, that \ingram{} allows a hybrid to maintain its loss with fewer
of them.

\subsection{State per session}
\label{sec:serving-kind}

\begin{figure}[tb]
\centering
\begin{tikzpicture}
\begin{groupplot}[
  group style={group size=2 by 1, horizontal sep=2.0cm},
  height=5.6cm,
  axis lines=left, axis line style={mute}, tick style={mute},
  ymajorgrids, grid style={black!8},
  tick label style={font=\footnotesize}, label style={font=\footnotesize},
  title style={font=\footnotesize},
  xmode=log, xmin=0.8, xmax=165,
  xtick={1.024,4.096,16.384,65.536,131.072}, xticklabels={1k,4k,16k,64k,128k},
  xlabel={Context length (tokens)},
]
\nextgroupplot[title={(a) Slots written per token of context}, width=0.36\textwidth,
  xtick={1.024,4.096,16.384,131.072}, xticklabels={1k,4k,16k,128k},
  ymin=0, ymax=6, ytick={0,2,4,6}, ylabel={Slots per token}]
\addplot[cat3!85!black, thick, mark=*, mark size=1.8pt] table[x=ktok, y=rate]{generated/state.dat};

\nextgroupplot[title={(b) State per session}, width=0.48\textwidth, clip=false,
  ymode=log, ymin=0.8, ymax=1500, ytick={1,10,100,1000}, yticklabels={1,10,100,1000},
  ylabel={State (MB)}]
\draw[black!25, densely dashed] (axis cs:16.384,0.8) -- (axis cs:16.384,1500);
\addplot[cat2, thick, mark=square*, mark size=1.8pt] table[x=ktok, y=h3]{generated/state.dat};
\addplot[cat3, thick, mark=triangle*, mark size=2.2pt] table[x=ktok, y=l7]{generated/state.dat};
\addplot[cat1, thick, mark=*, mark size=1.8pt] table[x=ktok, y=wp]{generated/state.dat};
\node[font=\footnotesize, anchor=west, text=cat2!85!black] at (axis cs:145,537) {3:1 hybrid};
\node[font=\footnotesize, anchor=west, text=cat3!70!black] at (axis cs:145,290) {7:1 + \ingram{}};
\node[font=\footnotesize, anchor=west, text=cat1!85!black] at (axis cs:145,45.5) {GDN + \ingram{}};
\end{groupplot}
\end{tikzpicture}

\caption{(a) Slots that \ingram{} has written per token of context, for held-out trajectories read
from their first token: 704 trajectories up to 16,384 tokens, 566 at 65,536 tokens and 134 at
131,072 tokens. The bound is six. (b) State per session of our three models with 8 blocks of width
256, in fp32, computed from these slot counts with the accounting of Table~\ref{tab:cost}. At
131,072 tokens it is 537~MB for the 3:1 hybrid, 314~MB for the 7:1 hybrid with \ingram{} and
46~MB for GDN with \ingram{}. The dashed line marks 16,384 tokens, the context at which the models were
trained and evaluated.}
\label{fig:state}
\end{figure}

A model with \ingram{} has two components beyond its other weights, a hashed table and a slot
store. The hashed table is a parameter: all sessions share it, and each token reads one row per cue length. The slot store
belongs to one session: it contains one vector for every cue the session has written, located
through a hash index. Each generated token reads and writes one slot per cue length, regardless of
the length of the context. An attention layer stores a key and a value for every earlier token,
and each generated token scores all of them, or an index over them when the attention is sparse.
The slot store therefore requires little memory bandwidth and can reside in ordinary system memory
or on storage, whereas the attention cache requires fast memory. The schedule of writes and reads
depends on token ids alone, and prefill therefore fills the slots in one pass over the prompt.

The slot store grows with the context, because a session adds a slot each time it encounters a new
cue, at most six per token. On held-out trajectories read from their first token, the measured
rate is below this bound and decreases with depth (Figure~\ref{fig:state}a). A session has
written 2.89 slots per token after 16,384 tokens, 1.96 after 65,536 tokens and 1.34 after 131,072
tokens, and between 65,536 and 131,072 tokens a new token adds 0.67 slots. After 131,072 tokens
the store therefore holds 1.34 vectors of width 64 per token of context, against a key and a value
of width 256 in one attention layer, a factor of 6.0; after 16,384 tokens the factor is 2.8. With
the state of the GDN blocks included, a session of 16,384 tokens occupies 12.7~MB for GDN with
\ingram{}, 46.2~MB for the 7:1 hybrid with \ingram{} and 67.5~MB for the 3:1 hybrid in fp32
(Appendix~\ref{sec:cost}), and a session of 131,072 tokens 46~MB, 314~MB and 537~MB
(Figure~\ref{fig:state}b). Slot vectors and table rows quantize to 4 bits at a cost of at most
0.0005 nats on seeds 0 and 1 (exploratory; Appendix~\ref{app:other}).

\subsection{Two scenarios}
\label{sec:serving-local}

\begin{table}[tb]
\centering
\footnotesize
\setlength{\tabcolsep}{4pt}
{\small (a) One user: largest context that fits in memory, in tokens}\\[3pt]
% Generated by scripts/serving_cost.py --tex (a roofline model; every input is an assumption). Do not edit.
\begin{tabular}{@{}lccc@{}}
\toprule
& 12 attention layers & 6 attention layers & 4 attention layers \\
Machine & (as released) & + \ingram{} & + \ingram{} \\
\midrule
RTX 4090, 24 GB + 128 GB RAM & 674k & 1.35M & 2.02M \\
Laptop, 128 GB unified memory & 434k & 744k & 1.06M \\
Mini PC, 128 GB unified memory & 631k & 1.10M & 1.56M \\
\bottomrule
\end{tabular}
\\[8pt]
{\small (b) Shared replica: state per session and sessions per replica}\\[3pt]
% Generated by scripts/serving_cost.py --tex (a roofline model; every input is an assumption). Do not edit.
\begin{tabular}{@{}lcccccc@{}}
\toprule
& \multicolumn{3}{c}{128k-token sessions} & \multicolumn{3}{c}{1M-token sessions} \\
\cmidrule(lr){2-4}\cmidrule(lr){5-7}
& GPU state & host state & sessions & GPU state & host state & sessions \\
Configuration & (GB) & (GB) & per replica & (GB) & (GB) & per replica \\
\midrule
24 attention layers (as released) & 4.5 & 0.0 & 558 & 32.7 & 0.0 & 76 \\
12 attention layers + \ingram{} & 2.5 & 0.4 & 988 & 16.7 & 3.3 & 149 \\
6 attention layers + \ingram{} & 1.6 & 0.4 & 1,605 & 8.6 & 3.3 & 289 \\
\bottomrule
\end{tabular}

\caption{Two scenarios from a roofline model. (a) A model shaped like Qwen3.8-Flash-Next (48
layers, 12 of them sparse attention as released) serving one user: weights and N-gram tables at
4.5 bits per parameter, attention cache in 16-bit floats, \ingram{} with slot width 512 in 16-bit
floats, 1.5 slots per token and 24 bytes of index per slot. The model may use 110~GB of the
laptop's memory and 115~GB of the mini PC's; the laptop has 546~GB/s of memory bandwidth and the
mini PC 273~GB/s. On the first machine the cache and the weights outside the experts are in 22~GB
of the GPU memory, and the experts, the tables and the slot store are in 120~GB of system memory.
(b) A model shaped like Kimi K3.1 Pro (69 linear and 24 sparse-attention layers as released) on a
replica of 32 accelerators with 141~GB of memory each, 90\% of it usable: weights of 1,561~GB as
released (4.5 bits per parameter on average), attention cache in 16-bit floats, sessions limited by the accelerator memory left after
the weights, removed attention layers become linear layers, and
\ingram{} uses slot width 1,024 in 16-bit floats with 1.5 slots per token and 24 bytes of index
per slot, stored in host memory.}
\label{tab:serving}
\end{table}

For one user on one machine, memory determines how much context fits (Table~\ref{tab:serving}a).
As released, the 12 attention layers of a model shaped like Qwen3.8-Flash-Next require 25.3~KB of
cache and index per token of context. With six attention layers the cache requires 12.7~KB per
token, and \ingram{} adds 1.6~KB per token and a table of 0.4~GB. On a laptop with 128~GB of
unified memory, of which the model may use 110~GB and the weights and $N$-gram tables occupy
99~GB, the largest context that fits increases from 434k tokens to 744k with six
attention layers and to 1.06M with four. Both scenarios assume 1.5 slots per token, above the 1.34
that we measure after 131,072 tokens.

The same accounting applies to a replica that serves many sessions (Table~\ref{tab:serving}b). For
a model shaped like Kimi K3.1 Pro, a session of 1M tokens occupies 32.7~GB on the accelerators,
32.3~GB of it attention cache and index. With half of the attention layers a session occupies
16.7~GB there and 3.3~GB of \ingram{} slots in host memory, and a replica accommodates 149
sessions instead of 76, or 289 with a quarter of the attention layers. A decode step
takes 12 to 14~ms in all three configurations at this context, because reading the weights
dominates it, and tokens per second therefore scale with the number of sessions.
